% !TEX root = main.tex

\documentclass[14pt,a4paper]{extarticle}
\usepackage[utf8]{inputenc}
\usepackage[T2A]{fontenc}
\usepackage[russian]{babel}
\usepackage{amsmath, amssymb}
\usepackage{geometry}
\usepackage{hyperref}
\usepackage{graphicx}
\usepackage{microtype}
\usepackage{listings}
\usepackage{xcolor}
\usepackage{setspace}
\usepackage{svg}
\usepackage{amsmath}
\usepackage{float}
\usepackage{makecell}
\usepackage{array} % Для вертикального выравнивания m{}
\usepackage{tcolorbox}
\usepackage{tabularx}
\newcolumntype{C}{>{\centering\arraybackslash}X}
\onehalfspacing


\geometry{top=2cm, bottom=2cm, left=1.2cm, right=1cm}
% Определение цветов для классической светлой темы
\definecolor{codegray}{rgb}{0.5,0.5,0.5}
\definecolor{backcolour}{rgb}{0.98,0.98,0.95} % Почти белый фон без желтизны
\definecolor{lightgray}{rgb}{0.95, 0.95, 0.95} % Светло-серый цвет в формате RGB

\lstset{
    language=C++,
    basicstyle=\ttfamily\small,     % Моноширинный шрифт
    keywordstyle=\color{blue},      % Синие ключевые слова
    commentstyle=\color{gray},      % Серые комментарии
    stringstyle=\color{red},        % Красные строки
    numbers=left,                   % Нумерация слева
    numberstyle=\tiny\color{gray},
    stepnumber=1,
    breaklines=true,                % Перенос длинных строк
    frame=none,                   % Рамка вокруг кода
    backgroundcolor=\color{lightgray},
    tabsize=4,
    showstringspaces=false
}

\tcbuselibrary{skins}